{\scriptsize\setlength{\tabcolsep}{3pt}
\begin{longtable}{rccccp{4.4cm}}
\caption{Casos de teste de \texttt{reg\_status.vwf}: visão geral.}
\label{tab:caso-reg-status}\\
\hline $t$ (ns) & \texttt{nrst} & \texttt{status\_reg\_in} & \texttt{status\_reg\_wr\_ena} & \texttt{status\_reg\_out} & Observação \\ \hline \endfirsthead
\hline $t$ (ns) & \texttt{nrst} & \texttt{status\_reg\_in} & \texttt{status\_reg\_wr\_ena} & \texttt{status\_reg\_out} & Observação \\ \hline \endhead
0--20 & 0 & 00000 & 00000 & 00000 &  \\
20--40 & 1 & 11111 & 11111 & 11111 & solta o reset; escreve 11111 (todos habilitados) \\
40--60 & 1 & 00000 & 00001 & 11110 & zera só RN \\
60--80 & 1 & 00000 & 00010 & 11100 & zera só RV \\
80--100 & 1 & 00000 & 00100 & 11000 & zera só RZ \\
100--120 & 1 & 00000 & 01000 & 10000 & zera só RH \\
120--140 & 1 & 00000 & 10000 & 00000 & zera só RC \\
140--160 & 1 & 10101 & 11111 & 10101 & escreve 10101 em todos \\
160--180 & 1 & 11111 & 01010 & 11111 & in = 11111, mas só RV e RH habilitados \\
180--205 & 1 & 00000 & 00000 & 11111 & ninguém habilitado: mantém \\
205--220 & 0 & 00000 & 00000 & 00000 & reset assíncrono \\
\hline
\end{longtable}}
\vspace{-0.9cm}
\fonte{Autoria própria. Saídas do fim de cada intervalo, conforme o modelo de referência e a simulação do Quartus II.}
